% SPDX-License-Identifier: MIT
\chapter{Commands, Menues und Tastaturkuerzel}\label{chap:commands}
\ChapterSummary{Hier lernen Sie, Benutzeraktionen von ihrer
Darstellung zu trennen. Ein Command kann von einem Button,
einem Menuepunkt und einem Shortcut verwendet werden,
ohne dass dreimal dieselbe Geschaeftslogik implementiert wird.}

\section{Aktionen als eigene Komponenten}
Ein \term{Command} besitzt einen Beschreibungstext,
einen Enabled-Zustand und einen synchronen
Ausführungs-Callback. Es ist keine visuelle
Schaltflaeche; sein Lebenszyklus gehört der
Anwendung, einem Besitzer oder einem
ausdruecklich gesteuerten Scope.

\begin{lstlisting}[language=SASDCpp,caption={Eine Aktion, mehrere Ausloeser}]
Command saveCommand{"Speichern"};
saveCommand.setOnExecuted([] {
    // Validieren und Speicherung ausloesen.
});
Window window;
auto& form = window.emplace<VBox>();
auto& save = form.emplace<Button>();
save.bindCommand(saveCommand);
\end{lstlisting}
Die Reihenfolge der lokalen Variablen ist
absichtlich so gewählt, dass das Fenster und
der gebundene Button vor dem Command
zerstört werden. Der Binding-Vertrag kennt
auch lebensdauersichere, nichtbesitzende
\term{Command::Reference}-Objekte; die
Anwendung sollte die Abhaengigkeiten
trotzdem einfach und gut nachvollziehbar halten.

\section{Enabled bedeutet auch semantisch gesperrt}
Mit \texttt{saveCommand.setEnabled(false)}
sperren Sie die Ausführung des Commands.
\texttt{execute()} liefert dann \texttt{false}.
Gebundene Buttons können Text- und Enabled-
Zustand aus dem Command übernehmen.
Ein solcher Vorgang ist verlasslicher, als nur
die visuelle Farbe eines Buttons zu ändern.

\begin{lstlisting}[language=SASDCpp]
saveCommand.setEnabled(formIsValid);
if (!saveCommand.execute()) {
    // Kommando war nicht ausfuehrbar.
}
\end{lstlisting}
Hier steht \texttt{formIsValid} für eine
bereits berechnete boolesche Variable der
Anwendung; die Validierung selbst liegt
ausserhalb des Commands.

\section{MenuModel und MenuBarModel}
Ein \term{MenuBarModel} besitzt Menues,
ein \term{MenuModel} besitzt Menueeintraege.
Ein Eintrag ist ein Command, ein Separator
oder ein rekursives Submenue. Die Modelle
sind semantisch: Der Core muss die Position
eines terminalbasierten Popups oder einer
grafischen Menueleiste nicht selbst kennen.

\begin{lstlisting}[language=SASDCpp,caption={Menuebaum aufbauen}]
MenuBarModel bar;
MenuModel& file = bar.appendMenu("Datei");
file.appendCommand(saveCommand);
file.appendSeparator();
MenuModel& advanced = file.appendSubmenu("Erweitert");
advanced.appendCommand(saveCommand);
\end{lstlisting}
Das Beispiel demonstriert die gemeinsame
Command-Identitaet; in einer echten
Anwendung waere ein zweimal angebotenes
Speichern eventuell weniger sinnvoll
als eine zweite Aktion. Menue-Hierarchien
werden im aktuellen Terminal-Host über
transiente Praesentationsansichten,
Hit-Tests und Controller umgesetzt.

\section{Anzeige eines Shortcuts ist nicht sein Routing}
Ein MenuItem kann ein Kürzel \emph{anzeigen}.
Das ist nicht dasselbe wie dessen
\emph{Registrierung}. Dafuer existiert
\term{ShortcutMap}; sie ordnet
\term{Shortcut}-Werte einem Command zu.
Der Host entscheidet, in welcher Phase
ein KeyEvent zuerst an Menue/Fokusziel
oder an eine globale ShortcutMap geht.
Diese Unterscheidung verhindert, dass
jedes sichtbare Kürzel automatisch
eine unsichtbare globale Eingaberoute erzeugt.

\begin{lstlisting}[language=SASDCpp]
ShortcutMap shortcuts;
const Shortcut saveKey{Key::f2, KeyModifier::none};
shortcuts.bind(saveKey, saveCommand);
file.appendCommand(saveCommand, saveKey);

// 'event' ist ein KeyEvent aus dem Host:
const bool dispatched = shortcuts.dispatch(event);
\end{lstlisting}
Dieser Ausschnitt zeigt die Kern-API;
\texttt{event} muss im Host als
\term{KeyEvent} vorliegen. Die Methode
\texttt{dispatch} fuehrt das zugeordnete
Command nur aus, wenn der KeyEvent wirklich
passt und das Command vorhanden und
aktiviert ist.

\section{Einmalige Ausführung sicherstellen}
Bei Tastenkombinationen und Menues besteht
sonst die Gefahr, dass eine einzige Eingabe
gleich zweimal zur Ausführung gelangt.
Ein Host sollte eine feste Routing-Reihenfolge
definieren und \term{handled/ignored}
respektieren. Bei einem offenen Menue
müssen Navigation, Escape, Untermenue-
Wechsel und Pointer-Gesten kohaerent
vorrangig behandelt werden.
Als belastbares Referenzbeispiel dient
\filepath{examples/terminal_form_demo.cpp}.

\section{Beobachtbare Command-Zustände}
Mit \texttt{observeState()} lassen sich Änderungen
am Text und am Enabled-Zustand synchron beobachten.
Die Rückgabe ist eine move-only
\term{StateSubscription}; bei ihrem
Destruktor wird die Beobachtung
automatisch getrennt. Solche RAII-Tokens
sind insbesondere für wiederverwendbare
Widgets wichtig, die nicht zu
dauerhaften Besitzern eines Commands
werden sollen.

\PracticalTip{Bei SASD Gatehold kann etwa
\emph{Regel anwenden} gleichzeitig Button,
Menuepunkt und Funktionstaste sein.
Definieren Sie die Prüfung der Berechtigung
und die eigentliche Aktion genau einmal
hinter dem Command, nicht je Ausloeser.
Ein UI-Command ersetzt aber keine serverseitige
Autorisierung.}
